% \chapter{Methodology}

% This is what I did to test and confirm my hypothesis.


% You may want to split this chapter into sub chapters depending on your design. I suggest you change
% the title to something more specific to your project.

% This is where you describe your design process in detail, from component/device selection to actual
% design implementation, to how you tested your system. Remember detail is important in technical
% writing. Do not just write I used a computer give the computer specifications or the oscilloscopes part
% number. Describe the system in enough detail so that someone else can replicate your design as well
% as your testing methodology.

% If you use or design code for your system, represent it as flow diagrams in text.

% \textbf{IMPORTANT:} Include a motivation for your selection of an appropriate technique, or engineering or IT tool to solve your project problem. Discuss any limitations if appropriate.
\chapter{Design}

\section{Kinematic Mapping}

\begin{algorithm}
\caption{Forward Reaching Algorithm}\label{alg:backward_reaching}
\begin{algorithmic}[1]
    \Require A sorted array $A$ of size $n$, and a target value $T$
    \Ensure The index of $T$ in $A$, or $-1$ if not found
    
    \Function{ForwardReaching}{$J$}
    \State $C\gets J.child$
        \If {$C$ is a hinge}
            \State $\vec{d}\gets C.\vec{pos}-J.\vec{pos}$
            \LineComment[1]{Project $\vec{d}$ onto the plane perpendicular to $C.\vec{axis}_{aux}$}
            \State $\vec{l}\gets\vec{d}-(\vec{d}\cdot C.\vec{axis}_{aux})C.\vec{axis}_{aux}$
            \State $mag\_l \gets |\vec{l}|$
            \LineComment[1]{Update joint orientation}
            \If {$mab\_l=0$}
                \State $J.\vec{axis}_{link}\gets C.\vec{axis}_{link}$
            \Else
                \State $J.\vec{axis}_{link}\gets\vec{l}/|\vec{l}|$
            \EndIf
            \State $J.\vec{axis}_{aux}\gets C.\vec{axis}_{aux}$
            \LineComment[1]{Check joint limits}
            \State $dir\gets\sgn((J.\vec{axis}_{link}\times C.\vec{axis}_{link})\cdot J.\vec{axis}_{aux})$
            \State $C.value\gets\cos^{-1}(J.\vec{axis}_{link}\cdot C.\vec{axis}_{link})dir$
            \If {$C$ is limited}
                \If {$C.value>C.high$}
                    \State $C.value\gets C.high$
                    \State $J.\vec{axis}_{link}\gets C.\vec{axis}_{link}\cos(C.value)$
                    \State $J.\vec{axis}_{link}\gets J.\vec{axis}_{link}+(C.\vec{axis}_{link}\times C.\vec{axis}_{aux})\sin(C.value)$
                \ElsIf {$C.value<C.low$}
                    \State $C.value\gets C.low$
                    \State $J.\vec{axis}_{link}\gets C.\vec{axis}_{link}\cos(C.value)$
                    \State $J.\vec{axis}_{link}\gets J.\vec{axis}_{link}+(C.\vec{axis}_{link}\times C.\vec{axis}_{aux})\sin(C.value)$
                \EndIf
            \Else
                \LineComment[1]{Argument angle}
                \State $C.value\gets\mod(C.value,360\degree)$
                \If {$C.value>180\degree$}
                    \State $C.value\gets C.value-360\degree$
                \EndIf
            \EndIf
            \LineComment[1]{Update joint position}
            \State $J.\vec{pos}\gets C.\vec{pos}-J.\vec{axis}_{link}J.len_{link}$
            \Statex

        \ElsIf {$C$ is a pivot}
            \State $J.\vec{axis}_{link}\gets C.\vec{axis}_{link}$
            \State $J.\vec{pos}\gets C.\vec{pos}-J.\vec{axis}_{link}J.len_{link}$
            \Statex

        \ElsIf {$C$ is prismatic}
            \State $J.\vec{axis}_{link}\gets C.\vec{axis}_{link}$
            \State $\vec{d}\gets C.\vec{pos}-J.\vec{pos}$
            \State $C.value\gets\vec{d}\cdot J.\vec{axis}_{link}$
            \LineComment[1]{Check joint limits}
            \If {$C$ is limited}
                \If {$C.value>C.high$}
                    \State $C.value\gets C.high$
                \ElsIf {$C.value<C.low$}
                    \State $C.value\gets C.low$
                \EndIf
            \EndIf
            \State $J.\vec{pos}\gets C.\vec{pos}-J.\vec{axis}_{link}C.value$
        \EndIf
        

    \EndFunction
\end{algorithmic}
\end{algorithm}
